\chapter{Conclusion}\label{ch:conclusion}
%% FUNCTION: Collects the recommendations of Chapters 4 and 5 in answer to the research question
%% and presents them in one decision table with a status column.

We asked how often Round Robin achieves EFX and under which conditions it is therefore a good
heuristic for EFX. Round Robin is attractive because people can carry it out by hand, and it
always guarantees EF1. It has, however, no constant guarantee for approximate EFX, so its value
as a heuristic for EFX can only be judged in the average case. We simulated $20{,}000$ instances
for each combination of two to five agents, up to four goods per agent, five levels of agreement
and six levels of dispersion, and varied one condition at a time around a reference cell with
four agents and twelve goods.

In the reference cell, Round Robin achieves EFX in $0.293$ of the instances. Whether this is a
good or a poor result depends on the conditions, and Table~\ref{tab:decision} collects them. Three
conditions favour Round Robin: little agreement among the agents, similar values together with a
number of goods divisible by the number of agents, and few goods per agent with distinct
favourite goods, where EFX is even guaranteed. Three conditions work against it: strong
agreement, strongly differing values and many agents. Under full agreement without
divisibility, Round Robin fails EFX with certainty, and reversing the order of the last round is
a remedy that is just as easy to carry out. Compared with dividing the goods somehow, Round Robin
is far ahead unless the agents agree strongly.

\begin{table}[ht]
  \centering
  \caption{When Round Robin is a good heuristic for EFX. Status: proven results hold for every
  valuation of the stated kind; measured results hold in our simulation design; exploratory
  results rest on rules fixed after individual cells were known.}
  \label{tab:decision}
  \small
  \begin{tabular}{p{4.3cm}p{5.6cm}p{2.2cm}l}
    \toprule
    Situation & Recommendation & Status & Section \\
    \midrule
    At most two goods per person, distinct favourite goods
      & Use Round Robin; EFX is guaranteed
      & proven & \ref{sec:bundlesize} \\
    Similar values, number of goods divisible
      & Round Robin or any split into equal bundle sizes is almost always EFX
      & proven (limit), measured & \ref{sec:dispersion-limit} \\
    Similar values, number of goods not divisible
      & Round Robin is better than a random split, but often narrowly misses EFX
      & measured & \ref{sec:divisibility}, \ref{sec:random} \\
    Strongly differing values
      & Divisibility no longer matters; Round Robin fails more often and more deeply
      & measured & \ref{sec:dispersion}, \ref{sec:divisibility} \\
    Little to moderate agreement
      & Use Round Robin; far better than a random split
      & measured & \ref{sec:agreement}, \ref{sec:random} \\
    Strong agreement
      & Expect Round Robin to fail more often and by more
      & measured & \ref{sec:agreement} \\
    Full agreement, number of goods not divisible
      & Never plain Round Robin; reverse the last round, which helps with few goods per person or similar values
      & proven (Round Robin), measured, exploratory (remedy) & \ref{sec:divisibility-identical}, \ref{sec:reversed} \\
    More goods per person
      & Raises the EFX rate, unless agreement is strong
      & measured, exploratory & \ref{sec:goods} \\
    More people, same goods per person
      & EFX becomes rarer; with agreement, failures also become deeper
      & measured & \ref{sec:agents} \\
    \bottomrule
  \end{tabular}
\end{table}

Our contributions are threefold. First, we measure how often Round Robin achieves EFX and how far
it misses it, for small numbers of agents and for correlated values, where existing results are
either asymptotic or concern other targets. Second, we prove three simple statements that delimit
the measurements: a guarantee for at most two goods per agent with distinct favourite goods, the
certain failure under identical valuations without divisibility, and the limit for vanishing
dispersion with equal bundle sizes. Third, we compare Round Robin with a random split and with a
reversed last round, which turns the average case results into advice for people who divide goods
by hand. All measured statements rest on one family of value distributions and, for four and five
agents, on the assumption that an EFX allocation exists (Section~\ref{sec:limitations}).
